\subsection{When a source point forces the cubic jet}
\label{sec:kt-source-to-cubic-jet}

A characteristic-zero point with nonzero opposite extension
entries does not, by itself, force the independent two-parameter
cubic jet of KT-ORDINARY-CUBIC-MASSEY-CONDITION.
The obstruction is present even in a flat domain with the exact
torus weights and involution signs.  A sufficient extraction
criterion is given afterwards, with its extra valuation and
coefficient-support hypotheses stated explicitly.

\begin{proposition}[KT-EQUIVARIANT-POINT-NOT-CUBIC-JET]
For every prime $p>3$, the complete local algebra
\[
 R=\mathbf Z_p[[x,y,z]]/(pz-x^2y)                         \tag{SJ.1}
\]
is a $\mathbf Z_p$-flat domain.  It has a continuous formal
$\mathbf G_m$ action induced by weights $(1,-1,1)$ on $(x,y,z)$,
algebraic on each quotient $R/\mathfrak m^n$,
an involution with signs $(+,-,-)$, and an augmentation to
$\mathbf Z_p$.  There is a continuous local characteristic-zero
specialization with $x,y$ both nonzero.

There is an equivariant quadratic jet
\[
 R\longrightarrow \mathbf F_p[s,t]/(s,t)^3,\qquad
             x\mapsto s,\quad y\mapsto t,\quad z\mapsto0.
                                                               \tag{SJ.2}
\]
There is no lift to $\mathbf F_p[s,t]/(s,t)^4$ with these
first-order values of $x,y$, even without requiring
equivariance of the lift.
\end{proposition}
\begin{proof}
The actions preserve the defining ideal: both terms of its
generator have weight one and involution sign minus.
The augmentation sets all three variables to zero.
The characteristic-zero point is
\[
                         x=p,\qquad y=p,\qquad z=p^2.
                                                               \tag{SJ.3}
\]
All three images are in the maximal ideal and satisfy (SJ.1),
so substitution defines a continuous local map to $\mathbf Z_p$.

For the domain assertion, filter the regular ambient ring
with positive weights
$\mathrm{wt}(p)=\mathrm{wt}(x)=\mathrm{wt}(y)=1$ and
$\mathrm{wt}(z)=2$.
This filtration is separated and cofinal with the
maximal-ideal filtration.  Its associated graded ring is
$\mathbf F_p[P,X,Y,Z]$.
Multiplicativity of initial forms gives
\[
 \operatorname{gr}R
      =\mathbf F_p[P,X,Y,Z]/(PZ-X^2Y).                  \tag{SJ.4}
\]
The polynomial on the right is primitive and degree one
in $Z$, hence irreducible by Gauss's lemma, since
$\gcd(P,X^2Y)=1$.  The polynomial ring is a unique
factorization domain, so this quotient is a domain.
The separated filtration then makes $R$ a domain.
The element $p$ survives, since its initial form $P$
is nonzero in (SJ.4).  Therefore $R$ is torsion-free,
and hence flat, over the discrete valuation ring
$\mathbf Z_p$.

Equation (SJ.2) is valid because its only relation has
degree three after reduction modulo $p$.
For a putative cubic lift, write
$x=s+O((s,t)^2)$ and $y=t+O((s,t)^2)$.
Its defining relation reduces to
\[
                      0=x^2y=s^2t\pmod{(s,t)^4},
                                                               \tag{SJ.5}
\]
which is false.  This proves the assertion.
\end{proof}

This example has no asserted Galois or automorphic realization.
Its role is exact: flatness, a characteristic-zero point,
the torus grading, the involution, and nonzero opposite
entries do not imply the desired independent cubic jet.
The term $pz$ can cancel $x^2y$ at the characteristic-zero
point and disappears in the characteristic-$p$ lifting problem.

\begin{proposition}[KT-SPARSE-VALUATIVE-CUBIC-EXTRACTION]
Retain the field, exponent and auxiliary-unit hypotheses
of KT-ORDINARY-CUBIC-MASSEY-CONDITION.
Let $\mathcal O$ be a finite coefficient DVR, with
uniformizer $\pi$, residue field $k$, and
$r_\pi=v_\pi(p)$.
Let $r:G_{\mathbf Q}\to
\mathrm{GL}_3(\mathcal O)\rtimes\{\pm1\}$ be continuous,
with right component $\eta$, determinant-one left component,
and residual representation $r_0$.
Assume its restriction to $G_E$ is unramified away from
$p$ and that it has the full ordered ordinary flag at $p$
with exact inertia characters $(\chi^e,1,\chi^{-e})$.
Here the flag is in the semidirect Borel sense: after
one local conjugation the left matrices for the full
rational decomposition group are upper triangular.

Suppose an integer $d\ge1$ satisfies
\[
                              r_\pi>3d,                       \tag{SJ.6}
\]
Assume $r(g)r_0(g)^{-1}\equiv1\pmod{\pi^d}$ for all $g$,
and, in the displayed global basis, the truncated logarithm of
$r(g)r_0(g)^{-1}$ modulo $\pi^{3d+1}$ has the sparse form
\[
\begin{split}
 L(g)={}&\pi^d\{u(g)X+v(g)Y\}
             +\pi^{2d}w(g)Z+\pi^{3d}W(g),\\
 &\hspace{30mm}W(g)\in\mathfrak{sl}_3(k).
\end{split}                                                   \tag{SJ.7}
\]
Here $r_0$ is lifted by the finite Teichm\"uller characters
before reduction, and (SJ.7) means there are no intermediate
coefficient powers below $\pi^{3d}$.
Then the cochains $u,v$ satisfy (CM.2), including the strict
and away-prime conditions on $v$, and
\[
                             \mathfrak o_3(u,v)=0.             \tag{SJ.8}
\]
In particular a point satisfying these explicit hypotheses
forces the graded, involutive two-parameter cubic jet, even
though it initially supplies only a one-parameter specialization.
\end{proposition}
\begin{proof}
The ring $\mathcal O/\pi^{3d+1}$ has characteristic $p$
by (SJ.6).  It contains a coefficient field $k$: the unique
lifts of residue elements satisfying $a^{|k|}=a$ form that
field, since the derivative is a unit and Frobenius is
additive in characteristic $p$.
Successive $\pi$-adic expansion therefore identifies it
with $k[\pi]/\pi^{3d+1}$.
The subalgebra generated by $\lambda=\pi^d$ is
$k[\lambda]/\lambda^4$.
The finite residual Teichm\"uller representation lies in
the constant coefficient field.
The truncated logarithm is the finite polynomial
$N-N^2/2+N^3/3$ for $N=r(g)r_0(g)^{-1}-1$.
The hypothesis modulo $\pi^d$ gives $N^4=0$ in this
quotient, so logarithms and exponentials terminate at
degree three.  No characteristic-zero analytic
logarithm or its convergence range is used.

Project the homomorphism identity first to the coefficient
of $\lambda$.  It gives the cocycle laws for $u$ and $v$.
Projection to $\lambda^2Z$ gives (CM.8).
At degree three put
\[
 m(g)=\tfrac12\{W(g)_{12}-W(g)_{23}\},\qquad
 n(g)=\tfrac12\{W(g)_{21}+W(g)_{32}\}.                  \tag{SJ.9}
\]
The $U,V$ lines are separately preserved by the residual
adjoint action.  The BCH calculation therefore gives
exactly (CM.11) for these $m,n$.
The other components of $W$ enter only linearly at this
order and can be discarded: none contributes to either
projected equation.  Substituting the independent monomials
$s,t,st,s^2t,st^2$ for the indicated five coefficient slots
now gives the required global homomorphism over $A_3$.
Its determinant, grading and involution follow from
the five trace-zero Lie vectors and (CM.3)--(CM.5).

It remains to verify the local conditions after this
separation of the parameters.
At every away-prime inertia subgroup $I_q\cap G_E$,
both $r$ and $r_0$ are the identity.  Every coefficient
in (SJ.7), including the projections (SJ.9), is therefore
zero there.  The extracted family is unramified over $E$.
For $v$, the larger rational inertia group has quotient
of order at most two after intersecting $G_E$, and its
coefficient character $\beta$ is unramified.
A homomorphism from that quotient into the additive
$p$-group $k$ is zero.  Hence $v$ is unramified on
the full rational inertia away from $p$ as required.

At $p$, choose an inertia element $\tau$ for which
$\psi(\tau),1,\psi(\tau)^{-1}$ are pairwise distinct.
This is possible because $(p-1)\nmid2e$.
The exact ordinary inertia characters reduce to these
constant characters in our characteristic-$p$ quotient.
The ordered eigenflag of $r(\tau)$ is therefore uniquely
defined in the neighborhood of the standard flag.
Its eigenprojectors are polynomials in $r(\tau)$ with
unit constant denominators.  Because (SJ.7) is defined
over $k[\lambda]/\lambda^4$, its normalized lower-unipotent
flag matrix is defined over the same subalgebra.
It is the flag furnished by ordinarity, so is stable
under the full decomposition group.

At first order the only negative Lie vector in (SJ.7)
is $Y$.  The flag equation consequently gives
$v(g)=\beta(g)a-a$ for one constant $a\in k$ on the full
decomposition group.  Conjugating by $\exp(\lambda aY)$
leaves
\[
 \lambda u(g)X+\lambda^2
       \{w(g)-\tfrac12a(1+\beta(g))u(g)\}Z
                       +\lambda^3W'(g).                       \tag{SJ.10}
\]
There is no remaining negative term through degree two.
Uniqueness of the normalized lower flag means that its
remaining correction starts at degree three.
The exact prescribed inertial diagonal implies that
the coefficient of $Z$ in (SJ.10) is zero on inertia;
thus $w$ has precisely the ordinary normalization used
in (CM.8).

The $V$ projection of $W'$ is a coboundary in
$k(\eta\psi^{-1})$: it is canceled by the degree-three
$V$ component of that lower flag matrix.
This projection is exactly the $st^2V$ coefficient of
the separated two-parameter family after its conjugation
by $\exp(t aY)$, since other components of $W$ cannot
interact before degree four.
A further degree-three $V$ conjugation makes the
extracted family upper triangular.
Its degree-three $U$ term is already upper triangular,
and its only diagonal correction is the $stZ$ term,
which is zero on inertia.  Thus its ordinary local
condition is exact.  This proves (SJ.8) by the cubic
criterion.
\end{proof}

\paragraph{What remains unproved for an arbitrary source.}
The actual source supplies a finite coefficient field,
the residual representation, and the local flags.
It does not currently supply (SJ.6), the sparse support
(SJ.7), or a theorem allowing other initial directions
to be removed while preserving the two specified
nonzero extension classes.
Passing to a ramified coefficient extension alone does
not establish (SJ.6): both $v_\pi(p)$ and the existing
entry valuations scale by the ramification degree.
The algebra (SJ.1) shows why these missing conditions
cannot be replaced just by flatness or by the existence
of a characteristic-zero point with nonzero edges.

Accordingly (SJ.8) is a proved conditional extraction
criterion, not an assertion that every remaining
primitive, nonessentially-self-dual TypeI point has
the independent cubic jet.
